\begin{solapartat}
\punts{0,2}
\begin{align*}
f_A &= \frac{1}{T_A} = \frac{1}{1} = 1\ \mathrm{Hz} \to \omega_A = 2\pi f_A = 2\pi\ \mathrm{rad/s}\\
f_\mathrm{B} &= 2f_A = 2\ \mathrm{Hz} \to \omega_B = 2\pi f_B = 4\pi\ \mathrm{rad/s}
\end{align*}

\punts{0,2} $\omega_A = \sqrt{\dfrac{k}{m_A}} \Rightarrow k = \omega_A^2 m_A = (2\pi)^2\times 0,1 = 4\pi^2\times 0,1 = 0,4\pi^2\ \mathrm{N/m}$

\punts{0,6} $\omega_B = \sqrt{\dfrac{k}{m_B}} \Rightarrow m_B = \dfrac{k}{\omega_B^2} = \dfrac{0,4\pi^2}{(4\pi)^2} = 0,025\ \mathrm{kg} = 25\ \mathrm{g}$

També és vàlid:

\punts{0,8} $\dfrac{\omega_A^2}{\omega_B^2} = \dfrac{m_B}{m_A} \Rightarrow m_B = m_A\dfrac{\omega_A^2}{\omega_B^2} = \dfrac{(2\pi)^2}{(4\pi)^2}\times 0,1 = 0,025\ \mathrm{kg} = 25\ \mathrm{g}$
\end{solapartat}

\begin{solapartat}
\punts{0,2}
\begin{align*}
v_{A,\max} &= A\omega_A = 0,05\times 2\pi = 0,31\ \mathrm{m/s}\\
v_{B,\max} &= A\omega_B = 0,05\times 4\pi = 0,63\ \mathrm{m/s}
\end{align*}

\punts{0,2} Si $f_\mathrm{B} = 2f_A \Rightarrow T_B = 0,5\ \mathrm{s}$
\[ T_A = 1\ \mathrm{s} \]

\punts{0,6}
\figura[0.8]{sol_fig1.png}

El gràfic es construeix atenent els valors màxims de la velocitat i als períodes de 1s i 0,5~s, respectivament. La diferència de fase és de $\pi$ radiants quan estan en oposició de fase. Això passa a 0,5~s i a 1,5~s (en aquests moments cada objecte està a un extrem diferent, amb $v=0$).
\end{solapartat}
